\section{Basic x86 Architecture and Assembly}

\subsection{ISA and Machine State}

\begin{definition}
    \textbf{ISA vs. Microarchitecture:} The \textbf{Instruction Set Architecture (ISA)} defines the programmer-visible contract between software and hardware (instructions, registers, memory model), independent of the underlying \textbf{microarchitecture} (hardware implementation like caches and pipelines).
\end{definition}

\begin{definition}
    \textbf{CISC vs. RISC:}
    \begin{itemize}
        \item \textbf{CISC (Complex Instruction Set Computer, e.g., x86):} Variable-length instructions (1--15 bytes), arithmetic instructions can access memory directly, and explicit hardware stack operations (\texttt{push}, \texttt{pop}).
        \item \textbf{RISC (Reduced Instruction Set Computer, e.g., ARM, RISC-V):} Fixed-size 32-bit instructions, load-store architecture (only load/store instructions access memory), and larger register files.
    \end{itemize}
\end{definition}

\begin{remark}
    \textbf{Evolution of x86:} Evolved with backward compatibility from 16-bit (\textbf{8086}) to 32-bit (\textbf{IA32}) and 64-bit (\textbf{x86-64}). Due to the 16-bit origin, Intel terminology defines a ``word'' as 16 bits (2 bytes).
\end{remark}

\begin{definition}
    \textbf{Programmer-Visible State (x86-64):}
    \begin{itemize}
        \item \textbf{Program Counter (\texttt{\%rip}):} 64-bit address of the next instruction to execute.
        \item \textbf{Register File:} 16 general-purpose 64-bit integer registers and vector registers (\texttt{\%xmm}, \texttt{\%ymm}) for floating-point/SIMD.
        \item \textbf{Condition Codes:} Single-bit status flags (\texttt{CF}, \texttt{ZF}, \texttt{SF}, \texttt{OF}) updated by arithmetic operations.
        \item \textbf{Memory:} Flat byte-addressable virtual address space storing code, data, and the stack.
    \end{itemize}
\end{definition}

\begin{remark}
    \textbf{Machine Data Types \& Operations:} Assembly has no aggregate types (arrays, structs). It operates only on primitive integer data (1, 2, 4, or 8 bytes), floating-point values, and untyped memory addresses. Instructions fall into three categories:
    \begin{enumerate}
        \item \textbf{Data movement} between registers and memory (\texttt{mov}, \texttt{push}, \texttt{pop}).
        \item \textbf{Arithmetic and logical operations} (\texttt{add}, \texttt{sub}, \texttt{imul}, shifts, bitwise ops).
        \item \textbf{Control transfer} (\texttt{jmp}, conditional branches \texttt{jX}, \texttt{call}, \texttt{ret}).
    \end{enumerate}
\end{remark}

\begin{remark}
    \textbf{AT\&T Syntax (GCC/Linux default):} Operands flow from source to destination:
    \[
        \texttt{instr}\langle\text{size}\rangle \quad \text{Source}, \quad \text{Destination}
    \]
    Registers use \texttt{\%} (\texttt{\%rax}), and immediates use \texttt{\$} (\texttt{\$42}). Size suffixes: \textbf{\texttt{b}} (1B, \texttt{char}), \textbf{\texttt{w}} (2B, \texttt{short}), \textbf{\texttt{l}} (4B, \texttt{int}), \textbf{\texttt{q}} (8B, \texttt{long}/pointer).
\end{remark}

\subsection{Registers and Data Movement}

\begin{definition}
    \textbf{x86-64 General-Purpose Registers:} Lower sub-registers (32-bit, 16-bit, 8-bit) can be accessed directly:
    \begin{center}
        \renewcommand{\arraystretch}{1.15}
        \setlength{\tabcolsep}{4pt}
        \begin{tabular}{|c|c|c|c|l|}
            \hline
            \textbf{64-bit (\texttt{q})} & \textbf{32-bit (\texttt{l})} & \textbf{16-bit (\texttt{w})} & \textbf{8-bit (\texttt{b})} & \textbf{Convention / Role} \\
            \hline
            \texttt{\%rax} & \texttt{\%eax} & \texttt{\%ax} & \texttt{\%al} & Return value \\
            \texttt{\%rbx} & \texttt{\%ebx} & \texttt{\%bx} & \texttt{\%bl} & Callee-saved \\
            \texttt{\%rcx} & \texttt{\%ecx} & \texttt{\%cx} & \texttt{\%cl} & 4th argument, shift count \\
            \texttt{\%rdx} & \texttt{\%edx} & \texttt{\%dx} & \texttt{\%dl} & 3rd argument \\
            \texttt{\%rsi} & \texttt{\%esi} & \texttt{\%si} & \texttt{\%sil} & 2nd argument \\
            \texttt{\%rdi} & \texttt{\%edi} & \texttt{\%di} & \texttt{\%dil} & 1st argument \\
            \texttt{\%rbp} & \texttt{\%ebp} & \texttt{\%bp} & \texttt{\%bpl} & Callee-saved (frame pointer) \\
            \texttt{\%rsp} & \texttt{\%esp} & \texttt{\%sp} & \texttt{\%spl} & Stack pointer \\
            \texttt{\%r8}--\texttt{\%r9} & \texttt{\%r8d}--\texttt{\%r9d} & \texttt{\%r8w}--\texttt{\%r9w} & \texttt{\%r8b}--\texttt{\%r9b} & 5th and 6th arguments \\
            \texttt{\%r10}--\texttt{\%r11} & \texttt{\%r10d}--\texttt{\%r11d} & \texttt{\%r10w}--\texttt{\%r11w} & \texttt{\%r10b}--\texttt{\%r11b} & Caller-saved scratch \\
            \texttt{\%r12}--\texttt{\%r15} & \texttt{\%r12d}--\texttt{\%r15d} & \texttt{\%r12w}--\texttt{\%r15w} & \texttt{\%r12b}--\texttt{\%r15b} & Callee-saved \\
            \hline
        \end{tabular}
    \end{center}
\end{definition}

\begin{important}
    \textbf{32-bit Zero-Extension:} Writing to a 32-bit register (e.g., \texttt{movl}) automatically zeroes the upper 32 bits of the 64-bit register. Writing to 8-bit or 16-bit registers preserves the upper bits.
\end{important}

\begin{definition}
    \textbf{Memory Addressing:} The general addressing mode $D(R_b, R_i, S)$ accesses memory at:
    \[
        \text{Mem}[\text{Reg}[R_b] + S \cdot \text{Reg}[R_i] + D]
    \]
    where $D$ is a constant displacement, $R_b$ is the base register, $R_i$ is the index register (cannot be \texttt{\%rsp}), and $S \in \{1, 2, 4, 8\}$ is the scale factor.
\end{definition}

\begin{definition}
    \textbf{Data Movement (\texttt{mov}):} Transfers data between immediates, registers, and memory (\texttt{movb}, \texttt{movw}, \texttt{movl}, \texttt{movq}). When moving a smaller value to a larger register:
    \begin{itemize}
        \item \textbf{\texttt{movz} (Zero-extension):} Fills upper bits with 0s (unsigned, e.g., \texttt{movzbl}).
        \item \textbf{\texttt{movs} (Sign-extension):} Replicates the sign bit into upper bits (signed, e.g., \texttt{movslq}, or \texttt{cltq} for \texttt{\%eax} $\to$ \texttt{\%rax}).
    \end{itemize}
\end{definition}

\begin{important}
    \textbf{No Memory-to-Memory Transfers:} A single \texttt{mov} instruction cannot have both source and destination in memory (\texttt{movq (\%rax), (\%rdx)} is illegal); use an intermediate register instead.
\end{important}

\subsection{Arithmetic and Logical Operations}

\begin{definition}
    \textbf{Load Effective Address (\texttt{leaq S, D}):} Uses memory addressing syntax to compute the address expression of \texttt{S} and stores the \emph{value itself} into register \texttt{D} ($D \leftarrow \&S$) \textbf{without accessing memory} or modifying condition codes. Used for pointer generation and fast linear arithmetic ($x + k \cdot y + c$).
\end{definition}

\begin{codeexample}
movq (%rdi, %rcx, 8), %rax  // Memory read: rax = rdi[rcx] (reads 8B from memory)
leaq (%rdi, %rcx, 8), %rax  // Address calc: rax = rdi + 8*rcx (no memory access)
leaq (%rdi, %rdi, 2), %rax  // Fast arithmetic: rax = 3 * rdi
\end{codeexample}

\begin{definition}
    \textbf{Standard Arithmetic \& Logical Instructions:}
    \begin{itemize}
        \item \textbf{Unary:} \texttt{incq D} ($D+1$), \texttt{decq D} ($D-1$), \texttt{negq D} ($-D$), \texttt{notq D} ($\sim D$).
        \item \textbf{Binary:} \texttt{addq S, D} ($D \leftarrow D + S$), \texttt{subq S, D} ($D \leftarrow D - S$), \texttt{imulq S, D} ($D \leftarrow D \times S$).
        \item \textbf{Logical:} \texttt{xorq}, \texttt{andq}, \texttt{orq}. (\texttt{xorl \%eax, \%eax} is the standard idiom to set \texttt{\%rax} $= 0$).
        \item \textbf{Shifts:} \texttt{salq}/\texttt{shlq} (left shift), \texttt{sarq} (arithmetic right shift, sign-extends), \texttt{shrq} (logical right shift, zero-fills). Variable shift counts must be in \texttt{\%cl}.
    \end{itemize}
\end{definition}

\subsection{Condition Codes, \texttt{setX}, and \texttt{jX}}

\begin{definition}
    \textbf{Condition Codes:} Single-bit flags set implicitly by arithmetic instructions (except \texttt{leaq}) or explicitly by comparison instructions:
    \begin{itemize}
        \item \textbf{CF (Carry Flag):} Unsigned overflow (carry out of MSB).
        \item \textbf{ZF (Zero Flag):} Result is zero ($== 0$).
        \item \textbf{SF (Sign Flag):} Result is negative ($\text{MSB} == 1$).
        \item \textbf{OF (Overflow Flag):} Signed two's-complement overflow.
    \end{itemize}
\end{definition}

\begin{definition}
    \textbf{Explicit Comparisons (\texttt{cmp} and \texttt{test}):} Update condition codes without altering operands:
    \begin{itemize}
        \item \textbf{\texttt{cmpq S2, S1}:} Sets flags based on $S_1 - S_2$ (note AT\&T operand order: compares $S_1$ with $S_2$).
        \item \textbf{\texttt{testq S2, S1}:} Sets flags based on $S_1 \ \& \ S_2$ (e.g., \texttt{testq \%rax, \%rax} checks if \texttt{\%rax} is $0$ or negative).
    \end{itemize}
\end{definition}

\begin{definition}
    \textbf{Reading Condition Codes (\texttt{setX}) and Branching (\texttt{jX}):}
    \begin{itemize}
        \item \textbf{\texttt{setX D}:} Sets a single low-order byte register (e.g., \texttt{\%al}) to \texttt{1} if condition \texttt{X} holds, or \texttt{0} otherwise (follow with \texttt{movzbl \%al, \%eax} to clear upper bits).
        \item \textbf{\texttt{jX Label}:} Conditionally jumps to \texttt{Label} if condition \texttt{X} holds (\texttt{jmp} for unconditional jump).
    \end{itemize}
    \begin{center}
        \renewcommand{\arraystretch}{1.15}
        \setlength{\tabcolsep}{6pt}
        \begin{tabular}{|l|l|l|l|}
            \hline
            \textbf{Suffix (\texttt{X})} & \textbf{Condition} & \textbf{Flags Checked} & \textbf{Description} \\
            \hline
            \texttt{e} / \texttt{z} & $==$ / Zero & $\text{ZF}$ & Equal / Zero \\
            \texttt{ne} / \texttt{nz} & $\neq$ / Non-zero & $\sim\text{ZF}$ & Not equal / Not zero \\
            \texttt{s} / \texttt{ns} & $< 0$ / $\ge 0$ & $\text{SF}$ / $\sim\text{SF}$ & Negative / Non-negative \\
            \hline
            \texttt{g} & Signed $>$ & $\sim(\text{SF} \oplus \text{OF}) \ \& \ \sim\text{ZF}$ & Greater (signed) \\
            \texttt{ge} & Signed $\ge$ & $\sim(\text{SF} \oplus \text{OF})$ & Greater or equal (signed) \\
            \texttt{l} & Signed $<$ & $\text{SF} \oplus \text{OF}$ & Less (signed) \\
            \texttt{le} & Signed $\le$ & $(\text{SF} \oplus \text{OF}) \mid \text{ZF}$ & Less or equal (signed) \\
            \hline
            \texttt{a} & Unsigned $>$ & $\sim\text{CF} \ \& \ \sim\text{ZF}$ & Above (unsigned) \\
            \texttt{ae} & Unsigned $\ge$ & $\sim\text{CF}$ & Above or equal (unsigned) \\
            \texttt{b} & Unsigned $<$ & $\text{CF}$ & Below (unsigned) \\
            \texttt{be} & Unsigned $\le$ & $\text{CF} \mid \text{ZF}$ & Below or equal (unsigned) \\
            \hline
        \end{tabular}
    \end{center}
\end{definition}
